\section{Reusing gCastle without changing the meaning of a discovered graph}
\label{exec20:gcastle-scope}

A bank comparing discovery tools needs to know what an output edge means before
it can decide whether a package belongs in the proposed platform. An adjacency
matrix is insufficient: the same two nonzero entries can represent an unresolved
orientation, a directed cycle, or a different graph convention. The earlier
Python survey identifies gCastle as a source of selected PC and gradient-based
comparators. The reuse choice developed here retains those method-specific
references and proposes a bounded original-engine experiment. It does not extend
the choice to the entire catalog or treat a discovered graph as an identified
intervention effect. The primary evidence is the gcastle subproject of
\texttt{huawei-noah/trustworthyAI} at its 29 May 2026 revision,
\texttt{fd8a04f3b57a0f00930406257d5c30dcb91da0a5}; the distribution is named
\texttt{gcastle}, while the import namespace is \texttt{castle}.
\citep{exec20_gcastle_source}

\subsection{Graph types and variable identity are part of the result}
\label{exec20:gcastle-encoding}

The inspected PC orientation routine documents a completed partially directed
acyclic graph (CPDAG). Under its encoding, $A_{ij}=1,A_{ji}=0$ represents
$i\to j$, while $A_{ij}=A_{ji}=1$ represents a single unresolved edge $i-j$.
The latter must neither be interpreted as two causal arrows nor converted to a
latent-confounding bidirected edge. A returned array also does not certify that
statistical assumptions or graph-class validity hold. The proposed result type
must therefore include the graph family, ordered variable identifiers,
orientation convention, restrictions and diagnostics alongside the matrix.
\citep{exec20_gcastle_source}

For an exact conversion reference, consider four illustrative variables:
offer $O$, response $R$, margin $M$ and cost shock $C$. In the ordering
$(O,R,M,C)$ let
\begin{equation}
 A=\begin{pmatrix}
 0&1&0&0\\
 1&0&1&0\\
 0&0&0&0\\
 0&0&1&0
 \end{pmatrix}.
 \label{exec20:gcastle-matrix}
\end{equation}
This represents $O-R\to M\leftarrow C$. Both choices of orientation for
$O-R$ produce acyclic completions with the displayed collider; neither is
selected by the unresolved pair alone. A decoder that treats every nonzero
entry as a directed edge instead reports the artificial cycle $O\to R\to O$.
Simultaneously permuting rows, columns and their labels to $(M,O,C,R)$ must
preserve all named edges. This property is checked by the authored reference;
it is not a report of native package conformance.
\citep{exec20_gcastle_reference}

The distinction matters when a bank combines several algorithms or asks an
agent to explain the diagram. A common graph object may hold unresolved edges,
but a downstream DAG-only estimator must receive a justified completion, an
effect set across admissible completions, or an explicit unsupported response.
It cannot silently select an orientation to satisfy an API. The inspected
\texttt{Tensor} constructor checks input container types and label lengths;
it does not explicitly validate unique labels, finite values or an exactly
two-dimensional input. Its documentation and finalization method also show that
operations such as slicing replace label metadata with numeric indices. The
wrapper must bind immutable identifiers before transformations and return the
explicit map with every result. Otherwise a correctly shaped matrix can be
attached to the wrong banking variables. \citep{exec20_gcastle_source}

\subsection{Support and knowledge restrictions require their own diagnostics}
\label{exec20:gcastle-support}

The PC skeleton implementation removes an edge when a conditional-independence
test returns $p\geq\alpha$. The original, stable and parallel variants are
separate choices; the stable path copies adjacency at each conditioning depth.
A supplied base skeleton is modified in place, including its diagonal, so a
service must copy it if the original request is to remain reproducible. Required
and forbidden prior edges also need validation before construction: the inspected
methods append to separate lists without an explicit contradiction or directed
cycle check. The orientation helper applies required entries and then forbidden
entries. These are concrete conversion obligations, not evidence that arbitrary
business knowledge has been translated into a consistent causal model.
\citep{exec20_gcastle_source}

The default Fisher-$z$ path computes a partial correlation using a matrix inverse
with a pseudoinverse fallback, clips the correlation, and multiplies its
transformation by $\sqrt{n-|Z|-3}$, where $n$ is the row count and $Z$ the
conditioning set. The inspected routine has no explicit guard requiring a
positive residual sample size. An integration contract must reject or separately
qualify constant variables, singular conditioning designs and insufficient
residual support, and must expose the chosen singularity policy. A pseudoinverse
is a numerical rule; it does not establish the reference distribution of the
test in a degenerate population. Dependence among repeated account histories
also needs a statistical treatment beyond these source-level conversions.
\citep{exec20_gcastle_source}

For categorical tests, the conditional power-divergence routine sums statistics
and degrees of freedom across observed strata, warns and skips strata raising
\texttt{ValueError}, and sets $p=1$ when the aggregate degrees of freedom are
zero. At $\alpha=0.05$, that value satisfies the PC edge-removal branch. This
literal source implication is useful for designing a prototype test, but is not
a native experimental result. It does not mean that the data supplied informative
evidence about independence throughout a target population. A bank segmentation
with almost no within-stratum variation needs an insufficient-support result
and a description of the population restriction. The service should retain
skipped strata, support counts and test diagnostics, rather than reducing all
such outcomes to a favorable independence verdict.
\citep{exec20_gcastle_source,exec20_gcastle_reference}

An LLM may propose a restriction such as ``the later offer cannot cause an earlier
balance.'' The proposal should retain the source of that temporal claim and
check it against the observation clock. It should reject contradictory
restrictions before invoking the engine. Natural-language confidence and the
existence of a prior-knowledge API provide neither the truth of the restriction
nor the validity of the conditional-independence test. The same typed restriction
and refusal must survive a direct library request, a CLI request and an MCP call.

\subsection{Numerical stopping and thresholding are different decisions}
\label{exec20:gcastle-numerics}

For the inspected linear NOTEARS routine, the fitted predictor is $XW$:
row $i$, column $j$ of $W$ gives the coefficient from variable $i$ to variable
$j$. The least-squares path centers the columns of $X$; logistic and Poisson
losses have separate formulas. The implemented acyclicity expression is the
polynomial
\begin{equation}
 h(W)=\operatorname{tr}\!\left[\left(I+
       \frac{W\circ W}{d}\right)^d\right]-d,
 \label{exec20:gcastle-polynomial}
\end{equation}
where $d$ is the number of variables and $\circ$ denotes elementwise
multiplication. The binary result uses the strict rule
$B_{ij}=\mathbf{1}\{|W_{ij}|>\tau\}$, with default $\tau=0.3$, and the weighted
matrix is retained separately. At equality the edge is absent, despite the
docstring's informal description of dropping weights below the threshold.
The inspected optimizer consumes \texttt{sol.x} without testing
\texttt{sol.success}; stopping may follow the residual tolerance, penalty bound
or iteration budget. The inspected routine does not separately validate final
binary topology. None of these observations establishes an actual failed fit.
\citep{exec20_gcastle_source}

An exact witness explains the extra topology requirement without assuming what
the optimizer returns. For two variables with off-diagonal weights $a=b=1/100$,
\begin{equation}
 h(W)=\frac{a^2b^2}{2}=\frac{1}{200000000}=5\times10^{-9}
 <10^{-8}.
 \label{exec20:gcastle-tolerance}
\end{equation}
At the explicitly chosen threshold $\tau=1/200$, both edges survive and form a
directed cycle. The default threshold $0.3$ removes both; this is not a
counterexample to default native behavior. It demonstrates that a positive
residual tolerance alone is not a topological certificate for arbitrary
threshold choices. The proposed worker must return the weighted result,
threshold, residual and actual topology status separately. If the public engine
surface does not expose required solver diagnostics, the service must report
that qualification is unavailable or use a separately reviewed instrumented
variant. A successful return cannot fill in missing convergence evidence.
\citep{exec20_gcastle_reference}

Units are equally consequential. Consider the stipulated linear equation
$Y=(3/10)X+\epsilon$, with $Y$ measured in dollars. A strict threshold of $3/10$
drops its coefficient. Expressing the predictor as $X'=X/100$ changes the same
equation to $Y=30X'+\epsilon$ and retains the coefficient at the same numerical
threshold. This is a statement about representation, not a claim that refitting
a penalized objective is scale invariant. Units, scaling and penalties must be
fixed together in a discovery comparison.
For an authored causal model in which the displayed arrow is indeed structural,
an intervention increment of 100 predictor units and an assumed cost of USD 20
give
\begin{equation}
 \Delta V=100(3/10)-20=\text{USD }10.
 \label{exec20:gcastle-value}
\end{equation}
Mistakenly transposing the graph assigns zero response in this simple model and
gives USD $-20$. This makes the economic consequence of a convention error
visible; it does not license treating an arbitrary fitted weight as an
identified effect. \citep{exec20_gcastle_reference}

\subsection{The reuse choice and a proportionate public prototype}
\label{exec20:gcastle-decision}

The selected source declares version 1.0.4, and its README dates that declaration
to 7 March 2025. The May 2026 monorepo commit does not by itself show recent
maintenance of each gCastle method. Environment declarations also differ:
the README states Python 3.6--3.9, while setup states only Python at least 3.6.
README and requirements include PyTorch; setup's installation dependencies omit
it, even though the inspected validator imports \texttt{torch} unconditionally.
The algorithm surface selects PyTorch by default and imports several gradient
methods. A proposed lightweight PC or NOTEARS worker therefore needs an actual
resolved and tested environment; a numerical routine's apparent simplicity
does not establish a minimal installation profile.
\citep{exec20_gcastle_source}

The original root license supplies Apache 2.0 terms, including grant and
redistribution conditions. Reuse must preserve applicable licenses, notices and
modification disclosures, and review the actual distribution and dependencies.
This source-level finding is not a completed rights inventory for all optional
backends. The catalog's static IID methods, event-sequence TTPM and HPCI marked
under development remain distinct evidence scopes.
\citep{exec20_gcastle_source}

The selected upstream test wrappers call a shared harness. Its PC and NOTEARS
methods iterate parameter combinations, catch exceptions and log them without
assertions or rethrow in those methods. That design is useful for exploring
execution problems, but an ordinary successful test-run status alone would not
establish graph correctness or even that all attempted fits succeeded. This
inspection makes no claim about every upstream test, and no downloaded test has
been run here. The proposed public prototype needs explicit expected outputs and
failures for graph decoding, variable permutations, threshold equality, units,
inconsistent priors, insufficient support and topology, plus separately designed
statistical comparisons on known generated mechanisms.
\citep{exec20_gcastle_source,exec20_gcastle_reference}

The bounded recommendation is to retain selected PC and linear NOTEARS as
references and evaluate a pinned original Python worker if those comparators
justify the prototype cost. Direct whole-catalog adoption would conceal
method-specific assumptions. Porting now would add responsibility for CI
degeneracy, orientation, stopping, thresholds and dependencies without a measured
platform advantage. A small independent exact reference remains valuable for
conversion checks, but is not a replacement implementation. The first prototype
can use public or generated observations; bank data is unnecessary to test these
software contracts. It should measure environment setup, serialization, support
diagnostics, optimization, memory and upgrade effort before any cost advantage
is asserted. Statistical usefulness and executable admission require their own
evidence after that prototype, while the present proposal can already state a
reasoned reuse choice.
